\section{Introduction}
\label{sec:intro}

Learning robot control from demonstrations turns a manipulation skill into supervised
learning: a policy is fit to the actions an operator took in the states the operator
visited~\citep{argall2009survey,pomerleau1991}. Two difficulties make this harder than it
looks. First, a policy that is slightly wrong drifts into states the demonstrations never
covered, and its errors compound over a closed-loop episode~\citep{ross2011dagger}. Second,
human demonstrations are often multimodal: several different action sequences solve the
same situation, and a regressor trained with a squared loss averages them into an action
that solves none~\citep{florence2021ibc,shafiullah2022bet}.

Diffusion policies~\citep{chi2023diffusion} address both concerns directly. They model the
conditional distribution of a whole chunk of future actions with a denoising diffusion
model~\citep{ho2020ddpm}, so a multimodal distribution can be represented, and they execute
the chunk in a receding-horizon loop, which reduces how often the policy must decide.
Reported gains over behavior cloning (BC) are large on several benchmarks, but a gain of
this kind can have more than one cause. A diffusion policy differs from a one-step BC
policy in its generative training objective, in its observation history, in predicting
action chunks rather than single actions, and often in its tuning budget. Action
chunking alone is known to help deterministic policies~\citep{zhao2023act}.

This report describes a controlled comparison that separates these factors as far as a
small budget allows. We compare four families on two robosuite tasks under one protocol:
a one-step MLP (BC-MLP), an LSTM over an observation window (BC-RNN), an MLP with exactly
the diffusion policy's observation and action horizons (BC-MLP-chunk), and a conditional 1D
temporal U-Net diffusion policy. The chunked MLP is the control that the protocol requires
before any advantage of diffusion over BC-MLP may be attributed to the generative model
(hypothesis H2-c). All families see the same exact-state observations, use the same action
encoding and parameter budget, and receive the same number of tuning trials; all are
selected by the same rule on validation rollouts and tested once on held-out initial
states after a committed freeze.

\paragraph{Contributions.} The contribution is empirical and methodological, not a new
algorithm:
\begin{itemize}
\item a preregistered, paired, closed-loop comparison of BC and diffusion policies on Lift
  and Can, with a chunked deterministic control, a held-out test run once, and every
  null result reported (Sections~\ref{sec:results}--\ref{sec:robustness});
\item preregistered robustness sets and ablations (exponential moving average, sampler,
  observation history), analysed with the same hierarchical paired bootstrap;
\item two separately frozen observation-modality extensions on Lift: a visual BC versus
  visual diffusion comparison, and a DP3-inspired point-cloud diffusion policy against
  the frozen RGB diffusion policy under a cube-placement shift
  (Section~\ref{sec:perception});
\item a reproducible pipeline: every number in this report is generated from hash-pinned
  frozen result files (Appendix~\ref{app:repro});
\item an independent audit of the code and frozen files after the results were frozen,
  whose findings are disclosed with their consequences (Section~\ref{sec:limitations}),
  including a seed-level sensitivity check of every declared success comparison, an analysis
  that was not preregistered and is labelled as such.
\end{itemize}
We make no claim of a new diffusion method, of real-robot performance, or of superiority
beyond what the preregistered paired intervals show.
